\subsection{酉表示}

定义 3. --- 设 G 是拓扑群，E 是 K 上的希尔伯特空间。G 在 E 中的酉表示，是 G 在 E 中的连续线性表示 \(\varrho\)，使得对 G 中每个 g，E 的自同态 \(\varrho(g)\) 都是 E 的酉自同态（EVT, V, p. 40）。

换言之，酉表示是希尔伯特空间中的等距表示。特别地，酉表示是有界的。

拓扑群在希尔伯特空间中的平凡表示是酉的。酉表示的每个子表示都是酉的。酉表示对子群的限制是酉的。

定义 4. --- 设 G 是拓扑群，\(\varrho\) 和 \(\pi\) 分别是 G 在希尔伯特空间 E 和 F 中的酉表示。E × F 上的 G-不变半双线性形式，是 E × F 上的连续半双线性形式 q，使得

\[
q (x, y) = q (\varrho (g) x, \pi (g) y)
\]

对所有 \(g \in \mathrm{G}\) 及每个 \((x, y) \in \mathrm{E} \times \mathrm{F}\) 成立。

E × F 上的 G-不变半双线性形式组成的向量空间记为 \(\mathrm{Sesq}_{\mathrm{G}}(\varrho,\pi)\) 或 \(\mathrm{Sesq}_{\mathrm{G}}(\mathrm{E},\mathrm{F})\)。

例. --- 设 G 是拓扑群，\(\varrho\) 是 G 在 E 中的酉表示。E 上的内积 \(q\) 是 \(\mathrm{E} \times \mathrm{E}\) 上的 G-不变半双线性形式。

引理 4. --- 设 G 是拓扑群，\(\varrho\) 是 G 到希尔伯特空间 E 的酉群 U(E) 中的同态。那么，\(\varrho\) 是 G 的酉表示，当且仅当对于 U(E) 上的逐点收敛拓扑，\(\varrho\) 在 G 的元素 e 处连续。只需该性质对 E 的一个总子集中的每个 \(x\) 成立即可。

该条件显然是必要的。根据 INT, VIII, p. 129, § 2, no. 1 的注 2，它也是充分的，因为 \(\varrho\) 的像在 \(\mathcal{L}(\mathrm{E})\) 中等度连续，并且映射 \(g \mapsto \varrho(g)x\) 在 \(e\) 处的连续性蕴含其在 G 上的连续性。

设 G 是拓扑群。如果 \(\varrho_{1}\) 和 \(\varrho_{2}\) 是 G 的酉表示，且 u 属于 \(\operatorname{Hom}_{\mathrm{G}}(\varrho_{1},\varrho_{2})\)，那么 \(u^{*}\) 属于 \(\operatorname{Hom}_{\mathrm{G}}(\varrho_{2},\varrho_{1})\)。事实上，由于 \(\varrho_{1}\) 和 \(\varrho_{2}\) 是酉表示，对每个 \(g\in G\) 有

\[
\begin{array}{r} u ^ {*} \circ \varrho_ {2} (g) = u ^ {*} \circ \varrho_ {2} (g ^ {- 1}) ^ {*} = (\varrho_ {2} (g ^ {- 1}) \circ u) ^ {*} \\ = (u \circ \varrho_ {1} (g ^ {- 1})) ^ {*} = \varrho_ {1} (g) \circ u ^ {*}. \end{array}
\]

引理 5. --- 设 G 是拓扑群，\(\varrho\) 是 G 在复希尔伯特空间 E 中的酉表示。空间 \(\mathrm{Hom}_{\mathrm{G}}(\varrho, \varrho)\) 是 \(\mathcal{L}(\mathrm{E})\) 的含单位星子代数。

前述内容表明，\(\operatorname{Hom}_{\mathrm{G}}(\varrho,\varrho)\) 是 \(\mathcal{L}(\mathrm{E})\) 的含单位自伴子代数。由于它在 \(\mathcal{L}(\mathrm{E})\) 中是闭的，所以它是 \(\mathcal{L}(\mathrm{E})\) 的星子代数。

引理 6. --- 设 \(\pi\) 和 \(\varrho\) 分别是拓扑群 G 在希尔伯特空间 E 和 F 中的酉表示。设 D 是 E 的稠密子空间，并且在 \(\pi\) 下稳定。设 u 是从 E 到 F、定义域为 D 的闭部分算子，并满足关系式 \(u \circ \pi(g) = \varrho(g) \circ u\) 对所有 \(g \in G\) 成立。则 \(u^{*}\) 的定义域在 \(\varrho\) 下稳定，并且有关系式 \(u^{*} \circ \varrho(g) = \pi(g) \circ u^{*}\) 对所有 \(g \in G\) 成立。

设 \(g \in \mathrm{G}\)，且 \(x \in \mathrm{dom}(u^*)\)。对于每个 \(y \in \mathrm{dom}(u)\)，有

\[
\begin{array}{r l} & {\langle \varrho (g) x \mid u (y) \rangle = \langle x \mid \varrho (g ^ {- 1}) u (y) \rangle = \langle x \mid u (\pi (g ^ {- 1}) y) \rangle} \\ & {\qquad = \langle u ^ {*} (x) \mid \pi (g ^ {- 1}) y \rangle = \langle \pi (g) (u ^ {*} (x)) \mid y \rangle ,} \end{array}
\]

因为 \(\varrho(g)^{*} = \varrho(g)^{-1} = \varrho(g^{-1})\)。这证明了 \(\varrho(g)x \in \mathrm{dom}(u^{*})\)，且 \(u^{*}(\varrho(g)x) = \pi(g)(u^{*}(x))\)。特别地，\(u^{*} \circ \varrho(g)\) 的定义域包含 \(u^{*}\) 的定义域，并且 \(\pi(g) \circ u^{*} \subset u^{*} \circ \varrho(g)\)。

此外，若 \(x \in \text{dom}(u^{*} \circ \varrho(g))\)，则 \(x = \varrho(g^{-1})(\varrho(g)x)\) 属于 \(\text{dom}(u^{*})\)，这是将前述结论应用于 \(g^{-1}\) 的结果。因此 \(u^{*} \circ \varrho(g) = \pi(g) \circ u^{*}\)。

命题 1. --- 设 \(\varrho_{1}\) 和 \(\varrho_{2}\) 分别是拓扑群 G 在希尔伯特空间 E \(_{1}\) 和 E \(_{2}\) 中的酉表示。从 Hom \(_{G}(\varrho_{1},\varrho_{2})\) 到 Sesq \(_{G}(\varrho_{2},\varrho_{1})\) 的映射将 u 关联到半双线性形式 \(q_{u}\)，该形式由 \(q_{u}(x,y)=\langle x|u(y)\rangle\) 定义；它是向量空间的同构。

若 \(u\) 是 G-态射，则有

\[
\begin{array}{c} q _ {u} (\varrho_ {2} (g) x, \varrho_ {1} (g) y) = \langle \varrho_ {2} (g) x \mid u (\varrho_ {1} (g)y) \rangle \\ = \langle \varrho_ {2} (g) x \mid \varrho_ {2} (g) u (y) \rangle = \langle x \mid u (y) \rangle = q _ {u} (x, y) \end{array}
\]

对所有 \(g \in G\) 及每个 \((x, y) \in \mathrm{E}_{2} \times \mathrm{E}_{1}\) 成立，因此所指映射是从 \(\operatorname{Hom}_{\mathrm{G}}(\varrho_{1}, \varrho_{2})\) 到 \(\operatorname{Sesq}_{\mathrm{G}}(\varrho_{2}, \varrho_{1})\) 的线性映射。由 EVT, V, p. 16, cor. 2，它是单射。

反之，设 \(q\) 是 G-不变半双线性形式，且 \(u\) 是从 \(\mathrm{E}_1\) 到 \(\mathrm{E}_2\) 的唯一线性映射，使得 \(q(x,y) = \langle x \mid u(y) \rangle\) 对所有 \((x,y) \in \mathrm{E}_2 \times \mathrm{E}_1\) 成立（同上）。对于每个 \(g\)（\(G\) 中）及每个 \((x,y)\)

在 \(\mathrm{E}_2\times \mathrm{E}_1\) 中，我们有

\[
\begin{array}{c} \langle x | (\varrho_ {2} (g) \circ u \circ \varrho_ {1} (g ^ {- 1})) (y) \rangle = \langle \varrho_ {2} (g ^ {- 1}) x | u (\varrho_ {1} (g ^ {- 1}) y) \rangle \\ = q (\varrho_ {2} (g ^ {- 1}) x, \varrho_ {1} (g ^ {- 1}) y) = q (x, y) = \langle x | u (y) \rangle , \end{array}
\]

因此 \(u = \varrho_{2}(g) \circ u \circ \varrho_{1}(g^{-1})\)。于是 \(u \in \operatorname{Hom}_{\mathrm{G}}(\varrho_{1}, \varrho_{2})\)。命题得证。

设 \(\varrho\) 是拓扑群 G 在希尔伯特空间 E 中的酉表示。若 K = C，记 \(\overline{\mathrm{E}}\) 为 E 的共轭空间（EVT, V, p. 6）。若 K = R，令 \(\overline{\mathrm{E}} = \mathrm{E}\)。共轭表示 \(\overline{\varrho}\) 是 G 在 \(\overline{\mathrm{E}}\) 中的表示，由 \(\overline{\varrho}(g) = \varrho(g)\)（对所有 \(g \in \mathrm{G}\)）定义。它是 G 的酉表示；按定义，\(\overline{\mathrm{E}}\) 的一个子空间是 \(\overline{\mathrm{E}}\) 的子表示，当且仅当它是 E 的子表示。

命题 2. --- 设 \(\varrho\) 是拓扑群 G 在希尔伯特空间 E 中的酉表示。记 u 为从 \(\overline{\mathrm{E}}\) 到 \(\mathrm{E}'\) 的等距同构，它将 \(x\) 关联到线性形式 \(y\mapsto \langle x|y\rangle\)。

\begin{enumerate}
\def\labelenumi{\alph{enumi})}
\item
  映射 u 是从共轭表示 \(\overline{\varrho}\) 到逆步表示 \(\breve{\varrho}\) 的同构；
\item
  给 E' 赋予通过 u 传递结构得到的希尔伯特空间结构；逆步表示 \(\breve{\varrho}\) 是 G 在 E' 中的酉表示。
\end{enumerate}

由 EVT, V, p. 15, th. 3 及其后注记，映射 \(u\) 是等距同构。

对 G 中每个 g、\(\overline{E}\) 中每个 x 以及 E 中每个 y，有

\[
\begin{array}{c} \langle y, (\breve {\varrho} (g) \circ u) (x) \rangle = \langle \varrho (g ^ {- 1}) y, u (x) \rangle \\ = \langle x | \varrho (g ^ {- 1}) y \rangle = \langle \varrho (g) x | y \rangle = \langle y, u (\overline {{\varrho}} (g) x) \rangle , \end{array}
\]

因此 \(\breve{\varrho}(g) \circ u = u \circ \overline{\varrho}(g)\)，这证明了 \(a\)；断言 \(b\) 立即得到。

推论. --- 设 \(\varrho\) 是拓扑群 G 在希尔伯特空间 E 中的酉表示。以下条件等价：

\begin{enumerate}
\def\labelenumi{(\roman{enumi})}
\item
  表示 \(\varrho\) 不可约；
\item
  逆步表示 \(\breve{\varrho}\) 不可约；
\item
  共轭表示 \(\overline{\varrho}\) 不可约。
\end{enumerate}

这由命题 2 以及命题前关于 \(\overline{E}\) 的子表示的注记推出。

命题 3. --- 设 \(\varrho\) 是拓扑群 G 在希尔伯特空间 E 中的酉表示。设 \(\pi\) 是 \(\varrho\) 的子表示，其表示空间为 F。F 在 E 中的正交补 F° 是 E 的子表示，且 E = F ⊕ F°，F° 与 E/F 同构。

空间 F° 是闭的。对于每个 \(x \in F^{\circ}\)、所有 \(g \in G\) 以及每个 \(y \in F\)，有 \(\langle \varrho(g)x \mid y \rangle = \langle x \mid \varrho(g^{-1})y \rangle = 0\)，因为 \(\varrho\) 是酉表示且 F 是子表示。因此 \(\varrho(g)x \in F^{\circ}\)。

E 到 E/F 的规范投影是 G-态射，所以它在 F° 上的限制是从 F° 到 E/F 的 G-态射；由 EVT, V, p. 13，它是等距同构。

我们也把 G 在 F° 中的表示记作 \(\pi^{\circ}\)。

命题 4. --- 设 \(\varrho\) 是拓扑群 G 在希尔伯特空间 E 中的酉表示。E 的闭子空间 F 是 \(\varrho\) 的子表示，当且仅当 E 上像为 F 的正交投影 p 是从 \(\varrho\) 到 \(\varrho\) 的 G-态射。

设 F 是 \(\varrho\) 的子表示。取 \(x \in E\)，写成 \(x = p(x) + y\)，其中 \(y \in F^{\circ}\)。对于 G 中每个 \(g\)，有

\[
\varrho (g) x = \varrho (g) (p (x)) + \varrho (g) y,
\]

并且由于 \(\varrho(g)(p(x))\) 属于 F 而 \(\varrho(g)y\) 属于 F°（命题 3），有 \(p(\varrho(g)x) = \varrho(g)(p(x))\)。因此 \(p\) 属于 \(\operatorname{Hom}_{\mathrm{G}}(\varrho, \varrho)\)。

反之，若 \(p \in \operatorname{Hom}_{\mathrm{G}}(\varrho, \varrho)\)，则 \(1_{E} - p\) 是 G-态射，因而 \(\mathrm{F} = \operatorname{Ker}(1_{\mathrm{E}} - p)\) 是 \(\varrho\) 的子表示。

